\documentclass[aspectratio=169,8pt]{beamer}
\usetheme{Madrid}
\usepackage[utf8]{inputenc}
\usepackage[T1]{fontenc}
\usepackage{amsmath,amssymb}
\usepackage{booktabs}
\usepackage{enumitem}
\setlist[enumerate]{label=\Alph*.,leftmargin=*,itemsep=1pt,topsep=2pt}
\setlist[itemize]{leftmargin=*,itemsep=1pt,topsep=2pt}
\setbeamerfont{block body}{size=\tiny}
\setbeamerfont{block title}{size=\scriptsize}
\setbeamerfont{frametitle}{size=\footnotesize}
\setbeamerfont{normal text}{size=\tiny}
\setbeamertemplate{navigation symbols}{}
\setbeamertemplate{itemize item}{\tiny$\bullet$}
\setbeamertemplate{itemize subitem}{\tiny$\circ$}

\title[GARP RAI]{GARP Risk \& AI --- Q451--Q550}
\subtitle{Unofficial Exam Prep}
\author{Unofficial}
\date{\today}

\begin{document}
	
	\begin{frame}
		\titlepage
	\end{frame}
	
	% =================================================================
	\section{Advanced Module 1 Topics}
	% =================================================================
	
	\begin{frame}{Q451 --- AI Risk Tolerance at Board Level}
		\begin{block}{Question}
			A board is defining its AI risk tolerance statement. Which of the following should be included as the most critical element?
			\begin{enumerate}
				\item A list of every AI model the firm owns.
				\item A clear statement of the types and levels of AI risk the organisation is willing to accept, aligned with strategy and regulatory obligations.
				\item Marketing language about the firm's AI capabilities.
				\item Employee satisfaction metrics.
			\end{enumerate}
		\end{block}
		\begin{exampleblock}{Answer: B}\end{exampleblock}
		\begin{block}{Explanation}
			AI risk tolerance defines acceptable risk levels and guides governance. Options A, C, and D are not the core of a risk tolerance statement.
		\end{block}
	\end{frame}
	
	\begin{frame}{Q452 --- Model Risk Appetite vs Tolerance}
		\begin{block}{Question}
			Which statement best distinguishes model risk appetite from model risk tolerance?
			\begin{enumerate}
				\item They are identical.
				\item Appetite defines the level of risk the organisation actively seeks or accepts to achieve objectives; tolerance defines the maximum deviation from appetite that triggers action.
				\item Appetite is only for technology risk.
				\item Tolerance is only for third-party risk.
			\end{enumerate}
		\end{block}
		\begin{exampleblock}{Answer: B}\end{exampleblock}
		\begin{block}{Explanation}
			Appetite = desired risk-return balance; tolerance = acceptable deviation. Options A, C, and D misstate the distinction.
		\end{block}
	\end{frame}
	
	\begin{frame}{Q453 --- Deep Learning and Financial Time Series}
		\begin{block}{Question}
			A team trains an LSTM on 20 years of hourly returns to forecast next-hour volatility. The model performs poorly out-of-sample. Which is the most likely reason?
			\begin{enumerate}
				\item LSTMs cannot handle time-series data.
				\item Financial returns are close to noise at hourly horizons, so complex models overfit non-repeating patterns.
				\item The dataset is too small.
				\item The model should be logistic regression.
			\end{enumerate}
		\end{block}
		\begin{exampleblock}{Answer: B}\end{exampleblock}
		\begin{block}{Explanation}
			Financial returns exhibit low signal-to-noise; complex models overfit. Options A, C, and D are incorrect or secondary.
		\end{block}
	\end{frame}
	
	\begin{frame}{Q454 --- Overfitting Detection in Deep Nets}
		\begin{block}{Question}
			A deep network achieves 99\% training accuracy and 62\% test accuracy. What action is most appropriate?
			\begin{enumerate}
				\item Increase model depth.
				\item Apply regularization (dropout, weight decay), reduce complexity, or obtain more data; use early stopping based on validation loss.
				\item Increase learning rate.
				\item Deploy immediately.
			\end{enumerate}
		\end{block}
		\begin{exampleblock}{Answer: B}\end{exampleblock}
		\begin{block}{Explanation}
			The large gap signals overfitting. Options A, C, and D would worsen it or ignore the problem.
		\end{block}
	\end{frame}
	
	\begin{frame}{Q455 --- Data Preprocessing Order}
		\begin{block}{Question}
			When preparing data for a distance-based model, which sequence is correct?
			\begin{enumerate}
				\item Scale, then split into train/test.
				\item Split into train/test, fit scaling on train, apply to test.
				\item Ignore scaling.
				\item Scale only test set.
			\end{enumerate}
		\end{block}
		\begin{exampleblock}{Answer: B}\end{exampleblock}
		\begin{block}{Explanation}
			Scaling parameters must be estimated on the training set only, then applied to the test set, to avoid leakage. Options A, C, and D are incorrect.
		\end{block}
	\end{frame}
	
	\begin{frame}{Q456 --- Target Leakage Detection}
		\begin{block}{Question}
			A model achieves suspiciously high test accuracy. Investigation reveals a feature that is computed using future information. What is the issue and remedy?
			\begin{enumerate}
				\item Overfitting; add regularization.
				\item Target leakage; the feature must be removed or recomputed using only information available at prediction time.
				\item Underfitting; add features.
				\item Data drift; retrain the model.
			\end{enumerate}
		\end{block}
		\begin{exampleblock}{Answer: B}\end{exampleblock}
		\begin{block}{Explanation}
			Future-information features cause leakage. Options A, C, and D do not address the root cause.
		\end{block}
	\end{frame}
	
	\begin{frame}{Q457 --- ML in Regulatory Reporting}
		\begin{block}{Question}
			A bank wants to use ML to compute regulatory capital. Which of the following is the most critical governance requirement?
			\begin{enumerate}
				\item Highest possible accuracy.
				\item Full compliance with regulatory standards: documentation, validation, interpretability, and independent review by qualified validators.
				\item Fastest training.
				\item Lowest cost.
			\end{enumerate}
		\end{block}
		\begin{exampleblock}{Answer: B}\end{exampleblock}
		\begin{block}{Explanation}
			Regulatory reporting requires compliance, documentation, validation, and interpretability. Options A, C, and D are insufficient.
		\end{block}
	\end{frame}
	
	\begin{frame}{Q458 --- Model Risk in Vendor Models}
		\begin{block}{Question}
			A firm uses a vendor-supplied AI model for credit scoring. Which governance gap is most likely?
			\begin{enumerate}
				\item The model cannot be tested.
				\item Limited transparency into model internals, restricting full independent validation; mitigation includes black-box testing and contractual rights to documentation.
				\item Vendor models cannot be used in regulated industries.
				\item Vendor models have no risk.
			\end{enumerate}
		\end{block}
		\begin{exampleblock}{Answer: B}\end{exampleblock}
		\begin{block}{Explanation}
			Vendor models limit transparency; black-box validation and contract terms help. Options A, C, and D are inaccurate.
		\end{block}
	\end{frame}
	
	\begin{frame}{Q459 --- Reproducibility in ML}
		\begin{block}{Question}
			Which of the following is essential for reproducibility of an ML model?
			\begin{enumerate}
				\item Only the model architecture.
				\item Fixed random seeds, documented data versions, environment specifications, and code versioning.
				\item Only training data.
				\item Only hyperparameters.
			\end{enumerate}
		\end{block}
		\begin{exampleblock}{Answer: B}\end{exampleblock}
		\begin{block}{Explanation}
			Reproducibility requires fixed seeds, data versions, environment, and code versioning. Options A, C, and D are partial.
		\end{block}
	\end{frame}
	
	\begin{frame}{Q460 --- Data Versioning}
		\begin{block}{Question}
			Why is data versioning important in AI governance?
			\begin{enumerate}
				\item To reduce storage.
				\item To enable reproducibility, auditability, and root-cause analysis of model behaviour over time.
				\item To speed up training.
				\item To reduce costs.
			\end{enumerate}
		\end{block}
		\begin{exampleblock}{Answer: B}\end{exampleblock}
		\begin{block}{Explanation}
			Data versioning supports reproducibility and audit. Options A, C, and D are secondary.
		\end{block}
	\end{frame}
	
	% =================================================================
	\section{Module 2 --- Advanced ML Techniques}
	% =================================================================
	
	\begin{frame}{Q461 --- K-Means Initialization}
		\begin{block}{Question}
			What is the purpose of the K-means++ initialization?
			\begin{enumerate}
				\item To randomly select centroids.
				\item To spread initial centroids far apart, reducing the chance of poor local optima and improving convergence.
				\item To eliminate outliers.
				\item To reduce K.
			\end{enumerate}
		\end{block}
		\begin{exampleblock}{Answer: B}\end{exampleblock}
		\begin{block}{Explanation}
			K-means++ chooses centroids with probability proportional to squared distance, spreading them out. Options A, C, and D misstate the purpose.
		\end{block}
	\end{frame}
	
	\begin{frame}{Q462 --- DBSCAN and Varying Densities}
		\begin{block}{Question}
			A dataset has clusters with very different densities. Which algorithm is likely to struggle?
			\begin{enumerate}
				\item K-means.
				\item DBSCAN, because it uses a single global density threshold.
				\item Hierarchical clustering.
				\item Gaussian Mixture Models.
			\end{enumerate}
		\end{block}
		\begin{exampleblock}{Answer: B}\end{exampleblock}
		\begin{block}{Explanation}
			DBSCAN struggles with varying densities because a single $\varepsilon$ cannot fit all clusters. Options A, C, and D are less affected.
		\end{block}
	\end{frame}
	
	\begin{frame}{Q463 --- Hierarchical Clustering Linkage Choice}
		\begin{block}{Question}
			Which linkage method is most sensitive to outliers?
			\begin{enumerate}
				\item Single linkage.
				\item Complete linkage.
				\item Average linkage.
				\item Ward's method.
			\end{enumerate}
		\end{block}
		\begin{exampleblock}{Answer: A}\end{exampleblock}
		\begin{block}{Explanation}
			Single linkage uses the closest points and can chain clusters via single points, including outliers. Options B, C, and D are more robust.
		\end{block}
	\end{frame}
	
	\begin{frame}{Q464 --- Elbow Method Limitations}
		\begin{block}{Question}
			What is a common limitation of the elbow method for choosing K?
			\begin{enumerate}
				\item It always produces a clear elbow.
				\item The elbow may be ambiguous or absent, and its identification is subjective.
				\item It only works for K=2.
				\item It requires labels.
			\end{enumerate}
		\end{block}
		\begin{exampleblock}{Answer: B}\end{exampleblock}
		\begin{block}{Explanation}
			Elbow plots may not have a clear bend. Options A, C, and D are inaccurate.
		\end{block}
	\end{frame}
	
	\begin{frame}{Q465 --- Silhouette Score vs Elbow}
		\begin{block}{Question}
			Why is the silhouette method often used alongside the elbow method?
			\begin{enumerate}
				\item Because both methods always agree.
				\item Because silhouette scores provide a quantitative criterion (higher is better) that complements the visual elbow heuristic.
				\item Because silhouette requires labels.
				\item Because elbow is never used.
			\end{enumerate}
		\end{block}
		\begin{exampleblock}{Answer: B}\end{exampleblock}
		\begin{block}{Explanation}
			Silhouette provides a numerical criterion to complement visual inspection. Options A, C, and D are inaccurate.
		\end{block}
	\end{frame}
	
	\begin{frame}{Q466 --- Random Forest Feature Importance}
		\begin{block}{Question}
			A random forest reports that a feature ``region\_id'' is highly important. Why should this be interpreted with caution?
			\begin{enumerate}
				\item Feature importance is always wrong.
				\item High-cardinality features can appear artificially important due to the impurity bias; permutation importance may give a more reliable ranking.
				\item Region has no business meaning.
				\item Only tree depth matters.
			\end{enumerate}
		\end{block}
		\begin{exampleblock}{Answer: B}\end{exampleblock}
		\begin{block}{Explanation}
			Impurity-based importance is biased toward high-cardinality features. Options A, C, and D are less accurate.
		\end{block}
	\end{frame}
	
	\begin{frame}{Q467 --- SVM Kernel Selection}
		\begin{block}{Question}
			Which kernel is most appropriate when the data is not linearly separable and the analyst has no prior knowledge of the boundary shape?
			\begin{enumerate}
				\item Linear kernel.
				\item RBF kernel.
				\item Polynomial of very high degree.
				\item No kernel.
			\end{enumerate}
		\end{block}
		\begin{exampleblock}{Answer: B}\end{exampleblock}
		\begin{block}{Explanation}
			The RBF kernel is a common default for non-linear separation. Options A, C, and D are less suitable defaults.
		\end{block}
	\end{frame}
	
	\begin{frame}{Q468 --- KNN Curse of Dimensionality}
		\begin{block}{Question}
			Which statement best describes the curse of dimensionality for KNN?
			\begin{enumerate}
				\item KNN gets faster as dimensionality increases.
				\item As dimensionality increases, distances become less discriminative and all points appear equidistant, degrading KNN performance.
				\item KNN is unaffected by dimensionality.
				\item KNN performs best in very high dimensions.
			\end{enumerate}
		\end{block}
		\begin{exampleblock}{Answer: B}\end{exampleblock}
		\begin{block}{Explanation}
			High dimensions cause distance concentration. Options A, C, and D are incorrect.
		\end{block}
	\end{frame}
	
	\begin{frame}{Q469 --- Regularization and Bias-Variance}
		\begin{block}{Question}
			Increasing the regularization strength in a linear model will:
			\begin{enumerate}
				\item Decrease bias and increase variance.
				\item Increase bias and decrease variance, potentially reducing overfitting.
				\item Have no effect.
				\item Increase both bias and variance.
			\end{enumerate}
		\end{block}
		\begin{exampleblock}{Answer: B}\end{exampleblock}
		\begin{block}{Explanation}
			Regularization trades bias for variance reduction. Options A, C, and D misstate the trade-off.
		\end{block}
	\end{frame}
	
	\begin{frame}{Q470 --- Cross-Validation for Hyperparameter Tuning}
		\begin{block}{Question}
			Why is cross-validation preferable to a single validation split for hyperparameter tuning?
			\begin{enumerate}
				\item It is faster.
				\item It reduces the variance of performance estimates and uses data more efficiently.
				\item It requires more data.
				\item It eliminates the need for a test set.
			\end{enumerate}
		\end{block}
		\begin{exampleblock}{Answer: B}\end{exampleblock}
		\begin{block}{Explanation}
			Cross-validation reduces variance of performance estimates. Options A, C, and D are incorrect or incomplete.
		\end{block}
	\end{frame}
	
	% =================================================================
	\section{Module 2 --- NLP and GenAI}
	% =================================================================
	
	\begin{frame}{Q471 --- Tokenization Errors}
		\begin{block}{Question}
			A pipeline fails to correctly identify negations. Which preprocessing step is most likely missing?
			\begin{enumerate}
				\item Tokenization.
				\item N-grams or negation handling --- bi-grams and above capture ``not good'' as a unit.
				\item Lowercasing.
				\item Feature extraction.
			\end{enumerate}
		\end{block}
		\begin{exampleblock}{Answer: B}\end{exampleblock}
		\begin{block}{Explanation}
			N-grams preserve local order and handle negation. Options A, C, and D do not address negation.
		\end{block}
	\end{frame}
	
	\begin{frame}{Q472 --- TF-IDF and Stop Words}
		\begin{block}{Question}
			A word appears in every document of a corpus. What is its TF-IDF score?
			\begin{enumerate}
				\item Very high.
				\item Zero (or near-zero), because IDF = log(N/N) = 0.
				\item Undefined.
				\item Equal to TF.
			\end{enumerate}
		\end{block}
		\begin{exampleblock}{Answer: B}\end{exampleblock}
		\begin{block}{Explanation}
			IDF = log(N/n) where n = document frequency. If n=N, IDF=0. Options A, C, and D are wrong.
		\end{block}
	\end{frame}
	
	\begin{frame}{Q473 --- Word Embeddings and Polysemy}
		\begin{block}{Question}
			A classic limitation of static word embeddings (Word2Vec) is:
			\begin{enumerate}
				\item They are too large.
				\item They assign a single vector per word, ignoring polysemy (multiple meanings).
				\item They cannot handle text.
				\item They always overfit.
			\end{enumerate}
		\end{block}
		\begin{exampleblock}{Answer: B}\end{exampleblock}
		\begin{block}{Explanation}
			Static embeddings assign one vector per word, missing context-dependent meanings. Options A, C, and D are incorrect.
		\end{block}
	\end{frame}
	
	\begin{frame}{Q474 --- Contextual Embeddings}
		\begin{block}{Question}
			Contextual embeddings from transformers differ from static embeddings by:
			\begin{enumerate}
				\item Having fewer dimensions.
				\item Producing different vectors for the same word depending on the surrounding context.
				\item Being faster.
				\item Using no training data.
			\end{enumerate}
		\end{block}
		\begin{exampleblock}{Answer: B}\end{exampleblock}
		\begin{block}{Explanation}
			Contextual embeddings depend on surrounding tokens; static embeddings do not. Options A, C, and D are incorrect.
		\end{block}
	\end{frame}
	
	\begin{frame}{Q475 --- RNN vs Transformer}
		\begin{block}{Question}
			Which best explains why transformers replaced RNNs in most NLP tasks?
			\begin{enumerate}
				\item Transformers use less memory.
				\item Transformers parallelise over sequence positions and capture long-range dependencies through attention.
				\item Transformers use fewer parameters.
				\item Transformers avoid training.
			\end{enumerate}
		\end{block}
		\begin{exampleblock}{Answer: B}\end{exampleblock}
		\begin{block}{Explanation}
			Transformers' parallelism and attention capabilities are the key advantages. Options A, C, and D are incorrect or secondary.
		\end{block}
	\end{frame}
	
	\begin{frame}{Q476 --- LLM Temperature}
		\begin{block}{Question}
			Setting the temperature of an LLM to a high value will:
			\begin{enumerate}
				\item Make outputs more deterministic.
				\item Increase randomness and creativity, at the cost of predictability.
				\item Decrease token count.
				\item Have no effect.
			\end{enumerate}
		\end{block}
		\begin{exampleblock}{Answer: B}\end{exampleblock}
		\begin{block}{Explanation}
			High temperature flattens the token distribution, increasing randomness. Options A, C, and D are incorrect.
		\end{block}
	\end{frame}
	
	\begin{frame}{Q477 --- Top-P Sampling}
		\begin{block}{Question}
			Top-P (nucleus) sampling works by:
			\begin{enumerate}
				\item Selecting all tokens with probability above a fixed threshold.
				\item Selecting the smallest set of tokens whose cumulative probability exceeds P.
				\item Selecting the K most probable tokens.
				\item Selecting tokens at random.
			\end{enumerate}
		\end{block}
		\begin{exampleblock}{Answer: B}\end{exampleblock}
		\begin{block}{Explanation}
			Top-P is cumulative-probability-based. Options A, C, and D describe other methods.
		\end{block}
	\end{frame}
	
	\begin{frame}{Q478 --- Hallucination Cause}
		\begin{block}{Question}
			Why do LLMs hallucinate?
			\begin{enumerate}
				\item They are trained on too little data.
				\item They optimise for statistical plausibility, not factual accuracy, generating fluent but incorrect outputs.
				\item They have no memory.
				\item They require more tokens.
			\end{enumerate}
		\end{block}
		\begin{exampleblock}{Answer: B}\end{exampleblock}
		\begin{block}{Explanation}
			Hallucinations stem from training objectives that optimise plausibility. Options A, C, and D are unrelated.
		\end{block}
	\end{frame}
	
	\begin{frame}{Q479 --- Retrieval-Augmented Generation}
		\begin{block}{Question}
			RAG improves LLM factual accuracy by:
			\begin{enumerate}
				\item Fine-tuning the model.
				\item Retrieving relevant documents at inference time and providing them as context.
				\item Increasing model size.
				\item Reducing context length.
			\end{enumerate}
		\end{block}
		\begin{exampleblock}{Answer: B}\end{exampleblock}
		\begin{block}{Explanation}
			RAG retrieves external knowledge and provides it to the model. Options A, C, and D do not describe RAG.
		\end{block}
	\end{frame}
	
	\begin{frame}{Q480 --- Prompt Engineering}
		\begin{block}{Question}
			Which is the most effective prompt engineering technique for structured output?
			\begin{enumerate}
				\item Short random prompts.
				\item Specifying role, format, and providing a few-shot example.
				\item Ignoring output format.
				\item Reducing context.
			\end{enumerate}
		\end{block}
		\begin{exampleblock}{Answer: B}\end{exampleblock}
		\begin{block}{Explanation}
			Role specification, format guidance, and examples reduce ambiguity. Options A, C, and D are ineffective.
		\end{block}
	\end{frame}
	
	\begin{frame}{Q481 --- Chain-of-Thought}
		\begin{block}{Question}
			Chain-of-thought prompting improves performance mainly because:
			\begin{enumerate}
				\item It changes model weights.
				\item It allows the model to use intermediate reasoning tokens to solve multi-step problems.
				\item It reduces computation.
				\item It increases randomness.
			\end{enumerate}
		\end{block}
		\begin{exampleblock}{Answer: B}\end{exampleblock}
		\begin{block}{Explanation}
			Chain-of-thought provides intermediate reasoning steps. Options A, C, and D are incorrect.
		\end{block}
	\end{frame}
	
	\begin{frame}{Q482 --- LLM Leaderboards}
		\begin{block}{Question}
			A key caveat of relying on LLM leaderboards is:
			\begin{enumerate}
				\item They are always definitive.
				\item Benchmark contamination and saturation can inflate or compress scores, and leaderboards may not reflect real-world performance.
				\item They ignore accuracy.
				\item They cover all tasks.
			\end{enumerate}
		\end{block}
		\begin{exampleblock}{Answer: B}\end{exampleblock}
		\begin{block}{Explanation}
			Benchmark issues include contamination, saturation, and lack of real-world relevance. Options A, C, and D are incorrect.
		\end{block}
	\end{frame}
	
	\begin{frame}{Q483 --- LLM Statelessness}
		\begin{block}{Question}
			What does ``stateless'' mean for an LLM?
			\begin{enumerate}
				\item It stores all previous conversations.
				\item It does not retain information between separate prompts unless prior context is re-supplied.
				\item It cannot process sequences.
				\item It works only offline.
			\end{enumerate}
		\end{block}
		\begin{exampleblock}{Answer: B}\end{exampleblock}
		\begin{block}{Explanation}
			LLMs process each prompt independently. Options A, C, and D are incorrect.
		\end{block}
	\end{frame}
	
	\begin{frame}{Q484 --- Agentic AI}
		\begin{block}{Question}
			What distinguishes an agentic AI system from a standard LLM application?
			\begin{enumerate}
				\item Higher accuracy.
				\item The ability to autonomously take actions in an environment with limited supervision.
				\item Lower cost.
				\item Smaller model size.
			\end{enumerate}
		\end{block}
		\begin{exampleblock}{Answer: B}\end{exampleblock}
		\begin{block}{Explanation}
			Agentic AI takes autonomous actions. Options A, C, and D are not the distinguishing feature.
		\end{block}
	\end{frame}
	
	\begin{frame}{Q485 --- GenAI Benchmarking}
		\begin{block}{Question}
			Why might a GenAI model score highly on a benchmark yet perform poorly in production?
			\begin{enumerate}
				\item Benchmarks are always accurate.
				\item Benchmarks may be contaminated, saturated, or unrepresentative of the real-world context in which the model is deployed.
				\item Benchmarks require no data.
				\item Benchmarks measure only speed.
			\end{enumerate}
		\end{block}
		\begin{exampleblock}{Answer: B}\end{exampleblock}
		\begin{block}{Explanation}
			Benchmarks have known limitations. Options A, C, and D are incorrect.
		\end{block}
	\end{frame}
	
	\begin{frame}{Q486 --- Fine-Tuning vs RAG}
		\begin{block}{Question}
			Which scenario best favours RAG over fine-tuning?
			\begin{enumerate}
				\item When proprietary knowledge changes frequently and must be updated often.
				\item When model weights must change permanently.
				\item When data is static.
				\item When latency is irrelevant.
			\end{enumerate}
		\end{block}
		\begin{exampleblock}{Answer: A}\end{exampleblock}
		\begin{block}{Explanation}
			RAG is best for frequently changing knowledge. Options B, C, and D are less aligned with RAG's strengths.
		\end{block}
	\end{frame}
	
	\begin{frame}{Q487 --- GenAI Bias Evaluation}
		\begin{block}{Question}
			Which is an appropriate method for evaluating bias in a GenAI system?
			\begin{enumerate}
				\item Ignore bias.
				\item Use demographic-parity-like metrics, stereotype scores, and adversarial prompts across groups.
				\item Only evaluate accuracy.
				\item Only evaluate speed.
			\end{enumerate}
		\end{block}
		\begin{exampleblock}{Answer: B}\end{exampleblock}
		\begin{block}{Explanation}
			Bias evaluation requires fairness metrics and adversarial testing. Options A, C, and D are insufficient.
		\end{block}
	\end{frame}
	
	\begin{frame}{Q488 --- GenAI Model Cards}
		\begin{block}{Question}
			A model card for a GenAI system typically includes:
			\begin{enumerate}
				\item Only the model's API key.
				\item Intended use, training data, evaluation results, limitations, and ethical considerations.
				\item Only training cost.
				\item Only temperature settings.
			\end{enumerate}
		\end{block}
		\begin{exampleblock}{Answer: B}\end{exampleblock}
		\begin{block}{Explanation}
			Model cards document intended use, training data, evaluations, limitations, and ethics. Options A, C, and D are partial.
		\end{block}
	\end{frame}
	
	\begin{frame}{Q489 --- GenAI Deployment Risks}
		\begin{block}{Question}
			Which of the following is a significant deployment risk for GenAI systems?
			\begin{enumerate}
				\item Uncontrolled output that could produce non-compliant or harmful content.
				\item Faster inference.
				\item Reduced storage.
				\item Lower cost.
			\end{enumerate}
		\end{block}
		\begin{exampleblock}{Answer: A}\end{exampleblock}
		\begin{block}{Explanation}
			Uncontrolled output is a key deployment risk. Options B, C, and D are benefits, not risks.
		\end{block}
	\end{frame}
	
	\begin{frame}{Q490 --- GenAI Prompt Injection}
		\begin{block}{Question}
			Prompt injection is:
			\begin{enumerate}
				\item A harmless input.
				\item A manipulation where an attacker inserts malicious instructions into prompts or data to cause unintended behaviour.
				\item A model architecture.
				\item A training technique.
			\end{enumerate}
		\end{block}
		\begin{exampleblock}{Answer: B}\end{exampleblock}
		\begin{block}{Explanation}
			Prompt injection attacks manipulate the model via crafted inputs. Options A, C, and D misstate.
		\end{block}
	\end{frame}
	
	% =================================================================
	\section{Module 3 --- Extended Risk Topics}
	% =================================================================
	
	\begin{frame}{Q491 --- Group vs Individual Fairness}
		\begin{block}{Question}
			Which is an example of individual fairness?
			\begin{enumerate}
				\item Equal approval rates across groups.
				\item Two applicants with identical qualifications receive the same decision.
				\item Equal TPR across groups.
				\item Equal FPR across groups.
			\end{enumerate}
		\end{block}
		\begin{exampleblock}{Answer: B}\end{exampleblock}
		\begin{block}{Explanation}
			Individual fairness = similar individuals treated similarly. Options A, C, and D are group fairness metrics.
		\end{block}
	\end{frame}
	
	\begin{frame}{Q492 --- Fairness Impossibility}
		\begin{block}{Question}
			The impossibility result for fairness holds:
			\begin{enumerate}
				\item Only when base rates are equal.
				\item When base rates differ, making it impossible to satisfy demographic parity, predictive rate parity, and equal opportunity simultaneously.
				\item Never.
				\item Only for regression models.
			\end{enumerate}
		\end{block}
		\begin{exampleblock}{Answer: B}\end{exampleblock}
		\begin{block}{Explanation}
			The impossibility arises under unequal base rates. Options A, C, and D are incorrect.
		\end{block}
	\end{frame}
	
	\begin{frame}{Q493 --- Data Composition Bias}
		\begin{block}{Question}
			Which is most likely to cause data composition bias?
			\begin{enumerate}
				\item Balanced dataset.
				\item Imbalanced dataset where minority class is underrepresented.
				\item Random data.
				\item Complete data.
			\end{enumerate}
		\end{block}
		\begin{exampleblock}{Answer: B}\end{exampleblock}
		\begin{block}{Explanation}
			Class imbalance drives composition bias. Options A, C, and D do not.
		\end{block}
	\end{frame}
	
	\begin{frame}{Q494 --- Measurement Bias}
		\begin{block}{Question}
			Which is an example of measurement bias?
			\begin{enumerate}
				\item Inconsistent measurement methods across groups.
				\item Random sampling.
				\item Balanced data.
				\item Complete data.
			\end{enumerate}
		\end{block}
		\begin{exampleblock}{Answer: A}\end{exampleblock}
		\begin{block}{Explanation}
			Measurement bias arises when features are measured differently across groups. Options B, C, and D are unrelated.
		\end{block}
	\end{frame}
	
	\begin{frame}{Q495 --- Historical Bias}
		\begin{block}{Question}
			Which is an example of historical bias?
			\begin{enumerate}
				\item Training on past hiring decisions that reflect discrimination against women.
				\item Random sampling.
				\item Balanced data.
				\item Normalized features.
			\end{enumerate}
		\end{block}
		\begin{exampleblock}{Answer: A}\end{exampleblock}
		\begin{block}{Explanation}
			Historical bias stems from past discrimination encoded in data. Options B, C, and D are unrelated.
		\end{block}
	\end{frame}
	
	\begin{frame}{Q496 --- Proxy Discrimination}
		\begin{block}{Question}
			Which is an example of proxy discrimination?
			\begin{enumerate}
				\item Explicit use of race as a feature.
				\item Use of a feature highly correlated with race (e.g., certain ZIP codes).
				\item Removing race.
				\item Balanced data.
			\end{enumerate}
		\end{block}
		\begin{exampleblock}{Answer: B}\end{exampleblock}
		\begin{block}{Explanation}
			Proxies carry discriminatory effects indirectly. Options A, C, and D are different or unrelated.
		\end{block}
	\end{frame}
	
	\begin{frame}{Q497 --- Interpretability vs Explainability}
		\begin{block}{Question}
			Which is true about a logistic regression model?
			\begin{enumerate}
				\item It is interpretable but often not explainable for complex cases.
				\item It is inherently interpretable and generally explainable.
				\item Neither interpretable nor explainable.
				\item Only explainable.
			\end{enumerate}
		\end{block}
		\begin{exampleblock}{Answer: B}\end{exampleblock}
		\begin{block}{Explanation}
			Logistic regression has interpretable coefficients and can be explained. Options A, C, and D are inaccurate.
		\end{block}
	\end{frame}
	
	\begin{frame}{Q498 --- Black-Box Models}
		\begin{block}{Question}
			Which technique is most commonly used to explain black-box model predictions?
			\begin{enumerate}
				\item Linear regression.
				\item SHAP or LIME.
				\item Logistic regression.
				\item Grid search.
			\end{enumerate}
		\end{block}
		\begin{exampleblock}{Answer: B}\end{exampleblock}
		\begin{block}{Explanation}
			SHAP and LIME are post-hoc explainability techniques. Options A, C, and D are unrelated.
		\end{block}
	\end{frame}
	
	\begin{frame}{Q499 --- Interpretability Trade-offs}
		\begin{block}{Question}
			The interpretability-accuracy trade-off implies:
			\begin{enumerate}
				\item Interpretable models are always less accurate.
				\item Complex models often achieve higher accuracy but lower interpretability; the trade-off should be managed based on context.
				\item Complex models are always better.
				\item Interpretability is irrelevant.
			\end{enumerate}
		\end{block}
		\begin{exampleblock}{Answer: B}\end{exampleblock}
		\begin{block}{Explanation}
			The trade-off is real and should be managed contextually. Options A, C, and D are oversimplifications.
		\end{block}
	\end{frame}
	
	\begin{frame}{Q500 --- Reputational Risk from AI}
		\begin{block}{Question}
			Which of the following is a direct source of AI-related reputational risk?
			\begin{enumerate}
				\item Increased storage.
				\item Privacy breaches, algorithmic bias, or failure to explain decisions.
				\item Faster training.
				\item Higher accuracy.
			\end{enumerate}
		\end{block}
		\begin{exampleblock}{Answer: B}\end{exampleblock}
		\begin{block}{Explanation}
			Reputational risk arises from privacy, bias, and explainability failures. Options A, C, and D are not sources of reputational risk.
		\end{block}
	\end{frame}
	
	\begin{frame}{Q501 --- Human Autonomy Violation}
		\begin{block}{Question}
			Which is a clear violation of human autonomy?
			\begin{enumerate}
				\item Explaining a decision.
				\item Manipulating a user's choices through dark patterns or AI-driven persuasion without consent.
				\item Providing choices.
				\item Showing accurate information.
			\end{enumerate}
		\end{block}
		\begin{exampleblock}{Answer: B}\end{exampleblock}
		\begin{block}{Explanation}
			Manipulative design undermines autonomy. Options A, C, and D respect autonomy.
		\end{block}
	\end{frame}
	
	\begin{frame}{Q502 --- Societal Risk Scenarios}
		\begin{block}{Question}
			Which of the following is a societal risk?
			\begin{enumerate}
				\item A bank loses money.
				\item An AI system amplifies political misinformation at scale.
				\item A chatbot mislabels a customer.
				\item A model underperforms.
			\end{enumerate}
		\end{block}
		\begin{exampleblock}{Answer: B}\end{exampleblock}
		\begin{block}{Explanation}
			Societal risk affects society at large. Options A, C, and D are individual or organisational risks.
		\end{block}
	\end{frame}
	
	\begin{frame}{Q503 --- Existential Risk}
		\begin{block}{Question}
			Which is a key concern in AI existential risk discussions?
			\begin{enumerate}
				\item Model speed.
				\item Superintelligent AI misaligned with human values.
				\item Storage costs.
				\item IT complexity.
			\end{enumerate}
		\end{block}
		\begin{exampleblock}{Answer: B}\end{exampleblock}
		\begin{block}{Explanation}
			Existential risk focuses on superintelligence and value alignment. Options A, C, and D are operational.
		\end{block}
	\end{frame}
	
	\begin{frame}{Q504 --- Global Inequality}
		\begin{block}{Question}
			AI may worsen global inequality because:
			\begin{enumerate}
				\item AI is universally accessible.
				\item Advanced AI capability concentrated in a few nations may widen economic and political gaps.
				\item AI reduces inequality.
				\item AI has no impact.
			\end{enumerate}
		\end{block}
		\begin{exampleblock}{Answer: B}\end{exampleblock}
		\begin{block}{Explanation}
			Concentration of AI can widen gaps. Options A, C, and D are inaccurate.
		\end{block}
	\end{frame}
	
	\begin{frame}{Q505 --- Misinformation Campaigns}
		\begin{block}{Question}
			Why is GenAI particularly concerning for misinformation campaigns?
			\begin{enumerate}
				\item It cannot generate text.
				\item It makes realistic content generation cheap and scalable.
				\item It only produces truthful content.
				\item It is too slow.
			\end{enumerate}
		\end{block}
		\begin{exampleblock}{Answer: B}\end{exampleblock}
		\begin{block}{Explanation}
			Generative AI makes content creation cheap and scalable. Options A, C, and D are incorrect.
		\end{block}
	\end{frame}
	
	\begin{frame}{Q506 --- Privacy and Surveillance}
		\begin{block}{Question}
			AI's ability to infer sensitive attributes from data creates which risk?
			\begin{enumerate}
				\item Faster computation.
				\item Privacy and surveillance risk.
				\item Lower storage costs.
				\item Higher accuracy.
			\end{enumerate}
		\end{block}
		\begin{exampleblock}{Answer: B}\end{exampleblock}
		\begin{block}{Explanation}
			Inference capabilities create privacy risks. Options A, C, and D are unrelated.
		\end{block}
	\end{frame}
	
	\begin{frame}{Q507 --- GenAI Risk Categories}
		\begin{block}{Question}
			Which of the following is a category of GenAI risk?
			\begin{enumerate}
				\item Functional risks.
				\item Data and information risks.
				\item Organizational and strategic risks.
				\item All of the above.
			\end{enumerate}
		\end{block}
		\begin{exampleblock}{Answer: D}\end{exampleblock}
		\begin{block}{Explanation}
			GenAI risks span functional, data, organizational, strategic, and societal categories. All three are included.
		\end{block}
	\end{frame}
	
	\begin{frame}{Q508 --- GenAI Data Risks}
		\begin{block}{Question}
			Which is a data risk unique to GenAI?
			\begin{enumerate}
				\item Storage cost.
				\item Training data may include proprietary, personal, or copyrighted material, creating legal exposure.
				\item Training time.
				\item Model size.
			\end{enumerate}
		\end{block}
		\begin{exampleblock}{Answer: B}\end{exampleblock}
		\begin{block}{Explanation}
			Copyright, privacy, and proprietary data in training corpora create legal exposure. Options A, C, and D are operational.
		\end{block}
	\end{frame}
	
	\begin{frame}{Q509 --- GenAI Organizational Risks}
		\begin{block}{Question}
			Which is an organizational risk of GenAI?
			\begin{enumerate}
				\item Responsibility gaps when multiple teams deploy GenAI without clear ownership.
				\item Faster inference.
				\item Lower cost.
				\item Reduced storage.
			\end{enumerate}
		\end{block}
		\begin{exampleblock}{Answer: A}\end{exampleblock}
		\begin{block}{Explanation}
			Responsibility gaps are a key organizational risk. Options B, C, and D are benefits.
		\end{block}
	\end{frame}
	
	\begin{frame}{Q510 --- GenAI Strategic Risks}
		\begin{block}{Question}
			Which is a strategic risk of GenAI?
			\begin{enumerate}
				\item Vendor lock-in and dependence on a small number of providers.
				\item Faster inference.
				\item Lower cost.
				\item Reduced storage.
			\end{enumerate}
		\end{block}
		\begin{exampleblock}{Answer: A}\end{exampleblock}
		\begin{block}{Explanation}
			Vendor lock-in is a strategic risk. Options B, C, and D are benefits.
		\end{block}
	\end{frame}
	
	% =================================================================
	\section{Module 4 --- Extended Ethics Topics}
	% =================================================================
	
	\begin{frame}{Q511 --- Utilitarianism}
		\begin{block}{Question}
			Which statement best captures utilitarianism in AI ethics?
			\begin{enumerate}
				\item It judges actions by adherence to rules.
				\item It judges actions by their consequences, maximising overall welfare.
				\item It focuses on character.
				\item It ignores outcomes.
			\end{enumerate}
		\end{block}
		\begin{exampleblock}{Answer: B}\end{exampleblock}
		\begin{block}{Explanation}
			Utilitarianism judges actions by consequences and aggregate welfare. Options A, C, and D are other frameworks or the opposite.
		\end{block}
	\end{frame}
	
	\begin{frame}{Q512 --- Deontology}
		\begin{block}{Question}
			Which statement best captures deontology in AI ethics?
			\begin{enumerate}
				\item Actions are judged by consequences.
				\item Actions are judged by adherence to duties and rules regardless of consequences.
				\item Focus on character.
				\item Focus on profit.
			\end{enumerate}
		\end{block}
		\begin{exampleblock}{Answer: B}\end{exampleblock}
		\begin{block}{Explanation}
			Deontology judges by duties and rules. Options A, C, and D are different frameworks.
		\end{block}
	\end{frame}
	
	\begin{frame}{Q513 --- Virtue Ethics}
		\begin{block}{Question}
			Which statement best captures virtue ethics?
			\begin{enumerate}
				\item Focus on consequences.
				\item Focus on duties.
				\item Focus on character traits and living a good life.
				\item Focus on laws.
			\end{enumerate}
		\end{block}
		\begin{exampleblock}{Answer: C}\end{exampleblock}
		\begin{block}{Explanation}
			Virtue ethics focuses on character. Options A, B, and D are different or incomplete.
		\end{block}
	\end{frame}
	
	\begin{frame}{Q514 --- Justice Conflicts}
		\begin{block}{Question}
			Why can different concepts of justice conflict?
			\begin{enumerate}
				\item They are identical.
				\item Egalitarianism, utilitarianism, and meritocracy prioritise different values, which can yield conflicting allocations.
				\item Only utilitarianism matters.
				\item Justice is unambiguous.
			\end{enumerate}
		\end{block}
		\begin{exampleblock}{Answer: B}\end{exampleblock}
		\begin{block}{Explanation}
			Different justice concepts conflict. Options A, C, and D misstate.
		\end{block}
	\end{frame}
	
	\begin{frame}{Q515 --- Autonomy and Consent}
		\begin{block}{Question}
			Which is essential for respecting autonomy in AI systems?
			\begin{enumerate}
				\item Ignoring user preferences.
				\item Obtaining informed consent and avoiding manipulative designs.
				\item Hiding decisions.
				\item Coercion.
			\end{enumerate}
		\end{block}
		\begin{exampleblock}{Answer: B}\end{exampleblock}
		\begin{block}{Explanation}
			Respecting autonomy requires consent and non-manipulation. Options A, C, and D violate autonomy.
		\end{block}
	\end{frame}
	
	\begin{frame}{Q516 --- Nonmaleficence}
		\begin{block}{Question}
			Which is an example of applying nonmaleficence in AI?
			\begin{enumerate}
				\item Deploying untested models.
				\item Testing AI thoroughly and removing unjustified risks before deployment.
				\item Ignoring safety.
				\item Racing to deploy.
			\end{enumerate}
		\end{block}
		\begin{exampleblock}{Answer: B}\end{exampleblock}
		\begin{block}{Explanation}
			Nonmaleficence requires avoiding harm, including unnecessary risk. Options A, C, and D violate the principle.
		\end{block}
	\end{frame}
	
	\begin{frame}{Q517 --- Beneficence}
		\begin{block}{Question}
			Which is an example of beneficence in AI?
			\begin{enumerate}
				\item Avoiding harm only.
				\item Actively designing AI to improve welfare, e.g., improving access to credit for underserved populations.
				\item Ignoring user needs.
				\item Maximising profit only.
			\end{enumerate}
		\end{block}
		\begin{exampleblock}{Answer: B}\end{exampleblock}
		\begin{block}{Explanation}
			Beneficence = actively doing good. Options A, C, and D are inconsistent with beneficence.
		\end{block}
	\end{frame}
	
	\begin{frame}{Q518 --- Privacy Best Practices}
		\begin{block}{Question}
			Which of the following is a privacy best practice?
			\begin{enumerate}
				\item Data minimisation and purpose limitation.
				\item Unlimited retention.
				\item Sharing with third parties without consent.
				\item Unencrypted storage.
			\end{enumerate}
		\end{block}
		\begin{exampleblock}{Answer: A}\end{exampleblock}
		\begin{block}{Explanation}
			Minimisation and purpose limitation are core privacy practices. Options B, C, and D violate them.
		\end{block}
	\end{frame}
	
	\begin{frame}{Q519 --- GDPR Data Subject Rights}
		\begin{block}{Question}
			Which is NOT a GDPR data subject right?
			\begin{enumerate}
				\item Access.
				\item Erasure.
				\item Portability.
				\item Right to sell data.
			\end{enumerate}
		\end{block}
		\begin{exampleblock}{Answer: D}\end{exampleblock}
		\begin{block}{Explanation}
			GDPR grants access, erasure, portability, rectification, restriction, and objection. Selling data is not a right.
		\end{block}
	\end{frame}
	
	\begin{frame}{Q520 --- EU AI Act Prohibited Practices}
		\begin{block}{Question}
			Which is explicitly prohibited under the EU AI Act?
			\begin{enumerate}
				\item Fraud detection.
				\item Social scoring and manipulative AI.
				\item Credit scoring.
				\item Customer service chatbots.
			\end{enumerate}
		\end{block}
		\begin{exampleblock}{Answer: B}\end{exampleblock}
		\begin{block}{Explanation}
			Social scoring and manipulative AI are prohibited. Other uses are regulated but not prohibited.
		\end{block}
	\end{frame}
	
	\begin{frame}{Q521 --- AI Ethics Framework Benefits}
		\begin{block}{Question}
			Which is NOT a benefit of an AI ethics framework?
			\begin{enumerate}
				\item Provides guidelines for ethical dilemmas.
				\item Builds trust.
				\item Reduces regulatory risk.
				\item Guarantees model accuracy.
			\end{enumerate}
		\end{block}
		\begin{exampleblock}{Answer: D}\end{exampleblock}
		\begin{block}{Explanation}
			Ethics frameworks do not guarantee model accuracy. Options A, B, and C are recognised benefits.
		\end{block}
	\end{frame}
	
	\begin{frame}{Q522 --- Ethical AI Committee Role}
		\begin{block}{Question}
			The primary purpose of an AI ethics committee is to:
			\begin{enumerate}
				\item Write code.
				\item Review high-risk use cases and provide guidance on ethical dilemmas.
				\item Market AI products.
				\item Train models.
			\end{enumerate}
		\end{block}
		\begin{exampleblock}{Answer: B}\end{exampleblock}
		\begin{block}{Explanation}
			Ethics committees review use cases and give guidance. Options A, C, and D are operational or unrelated.
		\end{block}
	\end{frame}
	
	\begin{frame}{Q523 --- Privacy Regulation Landscape}
		\begin{block}{Question}
			Which is an example of a US state-level privacy law?
			\begin{enumerate}
				\item GDPR.
				\item CCPA (California).
				\item PIPL (China).
				\item DSA (EU).
			\end{enumerate}
		\end{block}
		\begin{exampleblock}{Answer: B}\end{exampleblock}
		\begin{block}{Explanation}
			CCPA is California's privacy law. Options A, C, and D are non-US or different jurisdictions.
		\end{block}
	\end{frame}
	
	\begin{frame}{Q524 --- GDPR Fines}
		\begin{block}{Question}
			GDPR fines can reach:
			\begin{enumerate}
				\item €1,000.
				\item €100,000.
				\item Up to €20 million or 4\% of global annual turnover.
				\item No fines.
			\end{enumerate}
		\end{block}
		\begin{exampleblock}{Answer: C}\end{exampleblock}
		\begin{block}{Explanation}
			GDPR fines reach €20M or 4\% of global turnover. Options A, B, and D understate the maximum.
		\end{block}
	\end{frame}
	
	\begin{frame}{Q525 --- EU AI Act Fines}
		\begin{block}{Question}
			The EU AI Act's maximum fines reach:
			\begin{enumerate}
				\item €1 million.
				\item €5 million.
				\item Up to €35 million or 7\% of global turnover.
				\item No fines.
			\end{enumerate}
		\end{block}
		\begin{exampleblock}{Answer: C}\end{exampleblock}
		\begin{block}{Explanation}
			The AI Act allows fines up to €35M or 7\% of global turnover. Options A, B, and D are incorrect.
		\end{block}
	\end{frame}
	
	\begin{frame}{Q526 --- Responsible AI Principles}
		\begin{block}{Question}
			Which is a core Responsible AI principle?
			\begin{enumerate}
				\item Maximise profit only.
				\item Fairness, transparency, accountability, and privacy.
				\item Ignore stakeholders.
				\item Hide AI use.
			\end{enumerate}
		\end{block}
		\begin{exampleblock}{Answer: B}\end{exampleblock}
		\begin{block}{Explanation}
			Responsible AI includes fairness, transparency, accountability, and privacy. Options A, C, and D are contrary.
		\end{block}
	\end{frame}
	
	\begin{frame}{Q527 --- Local vs Global Explainability}
		\begin{block}{Question}
			Which is an example of global explainability?
			\begin{enumerate}
				\item SHAP for a single prediction.
				\item Feature importance across the full model.
				\item Counterfactual for one instance.
				\item LIME for one case.
			\end{enumerate}
		\end{block}
		\begin{exampleblock}{Answer: B}\end{exampleblock}
		\begin{block}{Explanation}
			Feature importance is a global technique. Options A, C, and D are local.
		\end{block}
	\end{frame}
	
	\begin{frame}{Q528 --- Explainability Requirements}
		\begin{block}{Question}
			In a regulated credit decision context, explainability requires:
			\begin{enumerate}
				\item Publishing all model internals.
				\item Meaningful justifications for individual decisions, typically via SHAP, counterfactuals, or adverse action notices.
				\item Only aggregate metrics.
				\item Refusing to explain.
			\end{enumerate}
		\end{block}
		\begin{exampleblock}{Answer: B}\end{exampleblock}
		\begin{block}{Explanation}
			Explainability requires meaningful per-decision justification. Options A, C, and D are insufficient or wrong.
		\end{block}
	\end{frame}
	
	\begin{frame}{Q529 --- Fairness Metrics Comparison}
		\begin{block}{Question}
			Which is a group fairness metric?
			\begin{enumerate}
				\item Individual fairness.
				\item Demographic parity.
				\item Consistency.
				\item Predictability.
			\end{enumerate}
		\end{block}
		\begin{exampleblock}{Answer: B}\end{exampleblock}
		\begin{block}{Explanation}
			Demographic parity is a group fairness metric. Options A, C, and D are different or unrelated.
		\end{block}
	\end{frame}
	
	\begin{frame}{Q530 --- Fairness Metric Conflict}
		\begin{block}{Question}
			Why can't all fairness metrics be satisfied simultaneously under unequal base rates?
			\begin{enumerate}
				\item Because fairness is subjective.
				\item Because mathematical constraints make demographic parity, predictive rate parity, and equal opportunity mutually incompatible.
				\item Because GDPR forbids it.
				\item Because of computational limits.
			\end{enumerate}
		\end{block}
		\begin{exampleblock}{Answer: B}\end{exampleblock}
		\begin{block}{Explanation}
			The impossibility result is mathematical. Options A, C, and D misstate the cause.
		\end{block}
	\end{frame}
	
	% =================================================================
	\section{Module 5 --- Extended Governance Topics}
	% =================================================================
	
	\begin{frame}{Q531 --- Data Governance Committee}
		\begin{block}{Question}
			Which best describes the data governance committee's role?
			\begin{enumerate}
				\item Operationally process data.
				\item Provide oversight on policies, standards, and practices, and resolve escalations.
				\item Build models.
				\item Market data.
			\end{enumerate}
		\end{block}
		\begin{exampleblock}{Answer: B}\end{exampleblock}
		\begin{block}{Explanation}
			The committee provides oversight and resolves escalations. Options A, C, and D are operational roles.
		\end{block}
	\end{frame}
	
	\begin{frame}{Q532 --- Data Quality Management}
		\begin{block}{Question}
			Which activity is part of data quality management?
			\begin{enumerate}
				\item Only storage.
				\item Data profiling, cleansing, validation, and monitoring.
				\item Only encryption.
				\item Only backup.
			\end{enumerate}
		\end{block}
		\begin{exampleblock}{Answer: B}\end{exampleblock}
		\begin{block}{Explanation}
			Data quality management includes profiling, cleansing, validation, and monitoring. Options A, C, and D are related but partial.
		\end{block}
	\end{frame}
	
	\begin{frame}{Q533 --- Data Lineage}
		\begin{block}{Question}
			Data lineage documents:
			\begin{enumerate}
				\item Only the file size.
				\item The origin, movement, and transformations of data across systems.
				\item Only the encryption status.
				\item Only the storage location.
			\end{enumerate}
		\end{block}
		\begin{exampleblock}{Answer: B}\end{exampleblock}
		\begin{block}{Explanation}
			Lineage tracks origin, movement, and transformations. Options A, C, and D are partial.
		\end{block}
	\end{frame}
	
	\begin{frame}{Q534 --- Metadata Repository}
		\begin{block}{Question}
			A metadata repository typically stores:
			\begin{enumerate}
				\item Raw data.
				\item Definitions, relationships, schemas, and technical/business metadata.
				\item Marketing content.
				\item Only encryption keys.
			\end{enumerate}
		\end{block}
		\begin{exampleblock}{Answer: B}\end{exampleblock}
		\begin{block}{Explanation}
			Metadata repositories store definitions, schemas, and relationships. Options A, C, and D are unrelated or partial.
		\end{block}
	\end{frame}
	
	\begin{frame}{Q535 --- Data Classification Levels}
		\begin{block}{Question}
			Which is a standard data classification level?
			\begin{enumerate}
				\item Public.
				\item Internal.
				\item Restricted.
				\item All of the above.
			\end{enumerate}
		\end{block}
		\begin{exampleblock}{Answer: D}\end{exampleblock}
		\begin{block}{Explanation}
			Standard classification: public, internal, confidential, restricted. All three are typical levels.
		\end{block}
	\end{frame}
	
	\begin{frame}{Q536 --- Data Access Control}
		\begin{block}{Question}
			The principle of least privilege requires:
			\begin{enumerate}
				\item Granting maximum access.
				\item Granting users only the access necessary for their role.
				\item No access.
				\item Full open access.
			\end{enumerate}
		\end{block}
		\begin{exampleblock}{Answer: B}\end{exampleblock}
		\begin{block}{Explanation}
			Least privilege = only what is needed. Options A, C, and D are inconsistent.
		\end{block}
	\end{frame}
	
	\begin{frame}{Q537 --- Data Encryption Standards}
		\begin{block}{Question}
			Which is a widely adopted standard for encrypting data at rest?
			\begin{enumerate}
				\item MD5.
				\item AES-256.
				\item SHA-1.
				\item Base64.
			\end{enumerate}
		\end{block}
		\begin{exampleblock}{Answer: B}\end{exampleblock}
		\begin{block}{Explanation}
			AES-256 is a standard symmetric encryption algorithm. Options A, C, and D are not encryption standards for confidentiality.
		\end{block}
	\end{frame}
	
	\begin{frame}{Q538 --- Model Inventory Best Practices}
		\begin{block}{Question}
			Which practice should be part of maintaining a model inventory?
			\begin{enumerate}
				\item Update the inventory only annually.
				\item Update on each material change, including new models, retirements, and changes in risk tier.
				\item Never update.
				\item Only capture model names.
			\end{enumerate}
		\end{block}
		\begin{exampleblock}{Answer: B}\end{exampleblock}
		\begin{block}{Explanation}
			The inventory is updated on material changes. Options A, C, and D are insufficient.
		\end{block}
	\end{frame}
	
	\begin{frame}{Q539 --- Model Risk Ranking}
		\begin{block}{Question}
			Which factor most raises the risk tier of an AI model?
			\begin{enumerate}
				\item Small user base.
				\item High materiality, high complexity, or high exposure.
				\item Short documentation.
				\item Low cost.
			\end{enumerate}
		\end{block}
		\begin{exampleblock}{Answer: B}\end{exampleblock}
		\begin{block}{Explanation}
			Risk tiering considers materiality, complexity, and exposure. Options A, C, and D do not raise risk.
		\end{block}
	\end{frame}
	
	\begin{frame}{Q540 --- Model Validation Scope}
		\begin{block}{Question}
			Which activity is part of independent model validation?
			\begin{enumerate}
				\item Rebuilding the model.
				\item Assessing conceptual soundness, data quality, implementation, and performance.
				\item Marketing the model.
				\item Entering data.
			\end{enumerate}
		\end{block}
		\begin{exampleblock}{Answer: B}\end{exampleblock}
		\begin{block}{Explanation}
			Validation covers conceptual soundness, data, implementation, and performance. Options A, C, and D are not validation activities.
		\end{block}
	\end{frame}
	
	\begin{frame}{Q541 --- Model Validation Report}
		\begin{block}{Question}
			A model validation report typically includes:
			\begin{enumerate}
				\item Validation opinion, findings, limitations, and recommendations.
				\item Only the model's accuracy.
				\item Only the code.
				\item Only the training data.
			\end{enumerate}
		\end{block}
		\begin{exampleblock}{Answer: A}\end{exampleblock}
		\begin{block}{Explanation}
			Validation reports include opinion, findings, limitations, and recommendations. Options B, C, and D are partial.
		\end{block}
	\end{frame}
	
	\begin{frame}{Q542 --- Model Risk Appetite Statement}
		\begin{block}{Question}
			Which is typically included in a model risk appetite statement?
			\begin{enumerate}
				\item A specific number of models.
				\item The firm's tolerance for model risk and escalation criteria.
				\item Marketing material.
				\item Only costs.
			\end{enumerate}
		\end{block}
		\begin{exampleblock}{Answer: B}\end{exampleblock}
		\begin{block}{Explanation}
			Appetite statements define tolerance and escalation criteria. Options A, C, and D are unrelated.
		\end{block}
	\end{frame}
	
	\begin{frame}{Q543 --- Model Risk Committee Membership}
		\begin{block}{Question}
			A model risk committee typically includes:
			\begin{enumerate}
				\item Only developers.
				\item Senior representatives from risk, business, validation, and possibly audit.
				\item Only IT.
				\item Only external consultants.
			\end{enumerate}
		\end{block}
		\begin{exampleblock}{Answer: B}\end{exampleblock}
		\begin{block}{Explanation}
			Cross-functional membership ensures diverse oversight. Options A, C, and D are incomplete.
		\end{block}
	\end{frame}
	
	\begin{frame}{Q544 --- Model Risk Reporting}
		\begin{block}{Question}
			Model risk reporting should be:
			\begin{enumerate}
				\item Occasional.
				\item Regular and ad hoc, with appropriate escalation to the board.
				\item Only annual.
				\item Only when problems occur.
			\end{enumerate}
		\end{block}
		\begin{exampleblock}{Answer: B}\end{exampleblock}
		\begin{block}{Explanation}
			Reporting is regular and ad hoc. Options A, C, and D are insufficient.
		\end{block}
	\end{frame}
	
	\begin{frame}{Q545 --- Model Risk Culture}
		\begin{block}{Question}
			Which is a hallmark of a strong model risk culture?
			\begin{enumerate}
				\item Silence about issues.
				\item Open communication, accountability, and continuous learning.
				\item Ignoring concerns.
				\item Only senior leaders caring.
			\end{enumerate}
		\end{block}
		\begin{exampleblock}{Answer: B}\end{exampleblock}
		\begin{block}{Explanation}
			A strong culture values openness, accountability, and learning. Options A, C, and D are contrary.
		\end{block}
	\end{frame}
	
	\begin{frame}{Q546 --- Model Risk Training}
		\begin{block}{Question}
			Model risk training should cover:
			\begin{enumerate}
				\item Only technical content.
				\item Technical, governance, ethical, and regulatory aspects.
				\item Only ethics.
				\item Only regulation.
			\end{enumerate}
		\end{block}
		\begin{exampleblock}{Answer: B}\end{exampleblock}
		\begin{block}{Explanation}
			Comprehensive training covers all four areas. Options A, C, and D are too narrow.
		\end{block}
	\end{frame}
	
	\begin{frame}{Q547 --- Change Management}
		\begin{block}{Question}
			A material model change should trigger:
			\begin{enumerate}
				\item Immediate deployment.
				\item Impact assessment, documentation, validation, and approval.
				\item No action.
				\item Retirement.
			\end{enumerate}
		\end{block}
		\begin{exampleblock}{Answer: B}\end{exampleblock}
		\begin{block}{Explanation}
			Material changes require structured review. Options A, C, and D violate governance.
		\end{block}
	\end{frame}
	
	\begin{frame}{Q548 --- Model Retirement}
		\begin{block}{Question}
			Model retirement requires:
			\begin{enumerate}
				\item Silent deletion.
				\item Formal documentation, inventory update, and archival.
				\item Continued use.
				\item No governance involvement.
			\end{enumerate}
		\end{block}
		\begin{exampleblock}{Answer: B}\end{exampleblock}
		\begin{block}{Explanation}
			Retirement is formal and documented. Options A, C, and D violate governance.
		\end{block}
	\end{frame}
	
	\begin{frame}{Q549 --- Responsible AI Governance}
		\begin{block}{Question}
			Responsible AI governance requires:
			\begin{enumerate}
				\item Only technical controls.
				\item A combination of technical controls, governance policies, and cultural elements.
				\item Only policies.
				\item Only culture.
			\end{enumerate}
		\end{block}
		\begin{exampleblock}{Answer: B}\end{exampleblock}
		\begin{block}{Explanation}
			Responsible AI governance integrates technical, policy, and cultural elements. Options A, C, and D are insufficient.
		\end{block}
	\end{frame}
	
	\begin{frame}{Q550 --- Closing Principle}
		\begin{block}{Question}
			The most important guiding principle in AI risk and governance is:
			\begin{enumerate}
				\item Maximise model complexity.
				\item Balance performance with governance, validation, monitoring, fairness, explainability, and ethical alignment with regulatory and business expectations.
				\item Ignore risks.
				\item Deploy all models without review.
			\end{enumerate}
		\end{block}
		\begin{exampleblock}{Answer: B}\end{exampleblock}
		\begin{block}{Explanation}
			Responsible AI balances performance and governance. Options A, C, and D contradict the principle.
		\end{block}
	\end{frame}
	
\end{document}